\chapter{Referencial Teórico}

\section{Ensino de Computação, inteligência artificial e formação na Educação Básica}

No ensino de Computação voltado à Educação Básica, o trabalho pedagógico difere do manuseio de ferramentas digitais porque investiga conhecimentos e práticas sobre o funcionamento de sistemas e de processos computacionais nas escolas. Esse debate já chega à rotina dos jovens. Quando os estudantes usam ferramentas automatizadas de inteligência artificial no dia a dia sem conhecer os princípios que orientam esses sistemas, essa convivência ainda difere do estudo próprio da Computação e da informática na educação analisada por \citeonline{valente2002uso}, em que o computador entra para apoiar outras disciplinas.

Na sequência didática planejada para o Ensino Médio, a plataforma interativa entra como material de apoio às aulas enquanto o estudo recai sobre o funcionamento introdutório dos modelos de \textit{machine learning} e sobre os limites desses sistemas. Operar uma tela não resolve esse problema. Com o uso de gráficos no plano cartesiano e funções ao lado de noções básicas de distância e erro médio nas atividades, a Matemática ajuda a turma a verificar como os dados e as regras alteram o resultado produzido pelo modelo.

Diretrizes curriculares recentes organizam essa área no país. Na regulamentação complementar à Base Nacional Comum Curricular que deixa a cargo das redes de ensino a forma de oferta em cada etapa escolar, a Resolução CNE/CEB nº 1/2022 estabelece normas nacionais para a Computação na Educação Básica \cite{brasil2022computacao}. Com foco na mediação do professor no trabalho escolar com sistemas inteligentes, o documento orientador publicado pelo Ministério da Educação em 2026 também afasta o uso da tecnologia como simples fornecedora de respostas prontas \cite{brasil2026iaeducacaobasica}. Essa base legal não prescreve algoritmos específicos. Segundo os referenciais da Sociedade Brasileira de Computação e da própria norma nacional, a área organiza-se pelos eixos Pensamento Computacional, Mundo Digital e Cultura Digital, cujas habilidades no Ensino Médio relacionam análise de dados e modelos computacionais à inteligência artificial \cite{sbc2018diretrizes,brasil2022computacao}.

Na formulação de problemas estudada por \citeonline{wing2006computational} no campo da Computação, o pensamento computacional participa desta proposta na organização da tarefa escolar pelos seus dados e regras por meio de práticas de abstração e decomposição. Esse recorte afasta leituras puramente técnicas. Quando separa uma situação estudada nos dados disponíveis e nas regras do modelo para compreender o conteúdo específico das aulas, o estudante mobiliza essas práticas sem que a pesquisa pretenda medir um desenvolvimento geral do pensamento computacional ou atestar ganhos diretos em outras disciplinas. No trabalho conjunto com os outros dois eixos da norma curricular durante os encontros na escola, o eixo Mundo Digital examina como os modelos processam representações numéricas e executam regras em decisões automatizadas e a Cultura Digital analisa criticamente os sistemas inteligentes presentes no cotidiano.

Na literatura internacional voltada à formação cidadã dos estudantes de acordo com cada etapa da Educação Básica, pesquisadores defendem que a escola discuta as ideias que orientam os sistemas de inteligência artificial e seus impactos na sociedade \cite{touretzky2019envisioning}. Esse debate inclui a aprendizagem a partir de dados. No documento da iniciativa AI4K12 dedicado ao ensino de \textit{machine learning} na escola, esse tema figura entre as grandes ideias da área e inclui para o Ensino Médio a diferenciação entre aprendizagem supervisionada e não supervisionada \cite{touretzky2023machine}. Tais referências internacionais não trazem um currículo pronto. Com a seleção de conteúdos introdutórios que dispensam programação avançada como requisito para as aulas, essas diretrizes somam-se às pesquisas brasileiras que mapearam aplicações na Educação Básica e relataram sequências didáticas sobre fundamentos da área no Ensino Médio \cite{camada2020ensino,duraes2021ensino}.

Ver uma mudança na tela não garante aprendizagem. Quando a observação visual do gráfico ocorre sem a tentativa de explicar a regra matemática e computacional em estudo, o avanço da turma passa a depender das previsões registradas na abertura da aula e da discussão conduzida pelo professor conforme o repertório matemático dos estudantes e a infraestrutura da escola. Quando discute em situações delimitadas como sistemas computacionais dependem de decisões humanas de modelagem e como alterações em parâmetros e regras afetam previsões e classificações, o estudo da inteligência artificial mantém a proposta vinculada ao ensino de Computação.

\section{\textit{Machine learning} introdutório no Ensino Médio}

O \textit{machine learning} encontra-se em um lugar específico entre os conteúdos de Inteligência Artificial, pois trata-se de examinar como os computadores são capazes de identificar padrões nos dados e utilizá-los para produzir respostas em tarefas definidas. \citeonline{marques2020teaching} realizaram um mapeamento de unidades instrucionais voltadas para a escola e constataram que existem vários materiais apenas com aspectos básicos de \textit{machine learning}, ou que apresentam de maneira muito abstrata o funcionamento do modelo. Por isso, no ensino dos fundamentos de \textit{machine learning}, discutir Inteligência Artificial significa permitir que os alunos possam compreender e visualizar um processo que muitas vezes só aparece pelo resultado final. Para que isso ocorra, esses sistemas recebem os dados de entrada e utilizam um modelo construído para devolver respostas dentro de um problema específico. A passagem de regras programadas para modelos orientados por dados é um desafio específico da educação em Computação, conforme o estudo de \citeonline{tedre2021teaching}, enquanto \citeonline{long2020ai} relacionam o letramento na área diretamente à capacidade de compreender e avaliar de forma crítica essas ferramentas.

A proposta do trabalho não envolve o detalhamento técnico de cada algoritmo, mas sim que o estudante consiga compreender o suficiente para vincular os dados disponíveis com a finalidade da tarefa, ao mesmo tempo em que verifica qual a regra de funcionamento do modelo. Outro ponto observado é que também se pretende que o aluno possa observar e verificar o efeito de ajustar os parâmetros ao interpretar a saída gerada. Dessa forma, a escolha desse recorte busca atender à dimensão voltada a compreender e avaliar o funcionamento dos sistemas de Inteligência Artificial discutida por \citeonline{ng2021conceptualizing}, sem que os estudantes precisem programar códigos complexos ou dominar técnicas avançadas de implementação.

De acordo com \citeonline{martins2023findings}, os resultados levantados na literatura mostram evidências de que o \textit{machine learning} pode ser trabalhado mesmo com estudantes iniciantes no Ensino Médio, quando eles compreendem os conceitos básicos e conseguem aplicá-los em tarefas relacionadas a algoritmos. \citeonline{sanusi2023systematic} realizaram uma revisão de experiências de implementação na Educação Básica e verificaram a presença expressiva de trabalhos voltados ao Ensino Médio. Outra revisão realizada por \citeonline{li2024systematic}, com 47 estudos empíricos, destaca o uso de tarefas práticas acompanhadas de apoio pedagógico como aspectos relevantes, mas também aponta limitações em relação aos recursos disponíveis para a execução das atividades e à profundidade das explicações sobre os conceitos. Muitas das pesquisas desenvolvidas na área apresentam falhas em comum, como a falta de inclusão do tema no currículo das escolas, a necessidade de preparo dos professores e a insuficiência na avaliação do impacto das aulas aplicadas. Para que a proposta possa ser bem-sucedida com a turma, independentemente da idade dos alunos, é importante: escolher bem o conteúdo; utilizar boas formas visuais de mostra dos dados; e estar atento para dar apoio durante a atividade. No Brasil, os autores \citeonline{martins2024low} realizaram um curso para 158 jovens de baixa renda, abrangendo o Ensino Médio. Mesmo com as diferenças para este projeto, os resultados mostram que é preciso adaptar as aulas à realidade local, ajudar os alunos individualmente, contar com a ajuda de outros para registrar como tudo ocorreu e que os cursos variam muito em termos de tempo e metas, e o que deu certo em uma escola, não funciona da mesma forma em outra, mas ainda assim a literatura mostra que a proposta tem boas condições de funcionar, cabendo à prática mostrar os resultados alcançados.

A sequência de aulas foi pensada para evitar iniciar com a exposição das fórmulas matemáticas dos algoritmos, pois sabe-se que ainda existem muitos alunos com dificuldades em relação a leitura de gráficos e trabalhar com proporções, além das limitações inerentes a aprendizagem por tela e o acúmulo de termos difíceis. Busca-se evitar que os alunos se desanimassem com as dificuldades iniciais e que a busca para apertar botões e tentar decorar termos difíceis escondessem a lógica por trás do modelo. As aulas iniciam trazendo situações próximas ao dia a dia dos estudantes, fazendo com que a turma tenha acesso aos dados e aos gráficos antes de tentar compreender a regra que produz a resposta final. Considera-se esse o início adequado para o Ensino Médio, sem deixar o ensino de Computação superficial. Além das dificuldades em encontrar os pontos nos dois eixos do gráfico, comparar distâncias e entender como as linhas e cores são utilizadas para mostrar a decisão do algoritmo, também é importante aprender a ler os gráficos. Aprender a ler os gráficos é parte da aula, pois apenas ver a tela mudar não garante que o aluno saiba explicar o motivo da alteração.

Outro desafio nas aulas é referente à base matemática dos alunos, pois muitas vezes as ideias de média, cálculo de distância, funções e proporção ainda causam dúvidas na turma e um curso curto de oito encontros não tem como recuperar todas as defasagens de anos anteriores. Por esse motivo, a proposta trabalha esses cálculos somente no momento em que são necessários para responder dúvidas práticas sobre o algoritmo, de maneira simples, sem transformar o percurso em um longo processo de revisão de Matemática separada da disciplina de Computação. Cada problema possui sua própria demanda matemática (como comparar o valor previsto com o valor real na regressão linear ou medir quão longe um ponto está de outro no \textit{k}-NN). O professor mediador avalia, na hora, se uma explicação rápida pode solucionar a dúvida dos alunos ou se a dificuldade trava o entendimento da atividade, como será detalhado na seção 3.3.

O uso da tela interativa, nesse sentido, ajuda a vencer alguns dos obstáculos citados na aprendizagem, porém pode esconder outros. Por exemplo, quando o aluno muda o controle para a quantidade de vizinhos, ele conseguirá ver o gráfico mudar na hora, mas a tela por si só não vai dizer o porquê de alguns pontos determinam a mudança do resultado final. Nesse ponto, será necessário a orientação do professor para explicar ao aluno o porquê da mudança (discutida mais detalhadamente na seção 3.3). A ferramenta desenvolvida reúne em si cinco diferentes técnicas (regressão linear, máquina de vetores de suporte -- \textit{SVM}, \textit{k}-vizinhos mais próximos -- \textit{k}-NN, árvore de decisão e \textit{k}-means) que possibilitam ao aluno trabalhar tanto a previsão de números e separação em classes quanto a formação de grupos em dados sem rótulo. Em todas as técnicas utilizadas, o aluno irá ler o problema, criar uma hipótese antes de testar na tela e posteriormente explicará como o fato de mexer nos dados ou ajustes feitos, alterando o gráfico. Os cinco algoritmos foram escolhidos e selecionados, pois funcionam de maneiras bem diferentes em gráficos fáceis de analisar, mesmo para os iniciantes. A regressão linear trabalha com a reta de tendência e a diferença entre valor previsto, o resíduo e o erro. A \textit{SVM} traça uma linha divisória entre dois grupos e mostra a margem até os pontos mais próximos. O \textit{k}-NN decide a classe ao olhar os exemplos mais parecidos com os arredores. A árvore de decisão segue uma sequência de perguntas e condições. E o \textit{k}-means junta os pontos em grupos, mesmo sem ter rótulos prévios.

A proposta consiste em trabalhar com dois modos diferentes de o computador aprender usando dados. Os primeiros quatro módulos (regressão linear, \textit{SVM}, \textit{k}-NN e árvore de decisão) trabalham com exemplos em que já possuem a resposta correta conhecida para que o modelo possa ser orientado para realizar a previsão ou classificação certa. Já o quinto módulo, com o uso do algoritmo \textit{k}-means, os dados não trazem nenhuma resposta ou classe pronta, antes do treinamento. Assim, torna-se necessário juntar os elementos em grupos, apenas ao medir a proximidade entre eles. De acordo com \citeonline{touretzky2023machine}, entender a diferença entre aprendizagem supervisionada e não supervisionada é um ponto fundamental no ensino de Inteligência Artificial. Dessa forma, nesta proposta, este projeto trabalhará a diferença entre as duas aprendizagens, por meio de conjuntos de dados sintéticos e recursos visuais.

Como já citado, os estudantes não vão aprender os cinco algoritmos no mesmo ritmo e do mesmo jeito. A regressão linear, por exemplo, é algo que remete a conteúdos de Matemática que eles já viram na escola, enquanto o agrupamento sem rótulos exige um pouco mais de tempo para que a turma perceba a diferença entre encontrar grupos novos e utilizar classes existentes anteriormente. Por isso, na aula seguinte aos exercícios, é importante revisar o assunto e, no sétimo encontro, comparar todos os modelos trabalhados para que cada técnica não pareça um assunto solto. Mesmo sendo voltada para iniciantes, as aulas mostram onde cada algoritmo pode falhar. Na regressão linear, por exemplo, os valores muito fora do padrão entortam a reta de tendência e é perigoso tentar prever valores longe da faixa de dados conhecida. Na \textit{SVM}, a linha que separa bem os pontos atuais da tela pode errar quando chegam dados novos. O resultado do \textit{k}-NN muda muito conforme os exemplos salvos, a escala usada nos eixos e a quantidade de vizinhos escolhida. Nas árvores de decisão, criar muitas regras pode atrapalhar em casos inéditos. E no \textit{k}-means, o formato dos grupos depende de quantos grupos foram pedidos e de onde os centroides começaram. Trabalhar esses limites durante a leitura dos gráficos ajuda a evitar tanto a ideia de que o computador acerta tudo sozinho quanto discussões apenas abstratas, pois ao mexer nos dados e ajustes de cada modelo o estudante vê na hora o efeito na tela e debate qual regra gerou aquele resultado.

\section{Representações, visualização e experimentação guiada}

O uso de recursos visuais facilita o primeiro contato dos alunos iniciantes com o aprendizado de máquina. \citeonline{gressevonwangenheim2021visual} realizaram um mapeamento de ferramentas voltadas à criação interativa de modelos, porém verificaram que a grande maioria delas se limita a classificar imagens e realizar tarefas rápidas. Diferente da maioria das ferramentas criadas, a desenvolvida neste projeto possui a particularidade de combinar conjuntos de dados sintéticos com tabelas montadas pela própria turma, para trabalhar regressão, classificação e agrupamento. Mesmo com o intuito de ampliar os horizontes desses temas, essa mudança sozinha não será suficiente para alcançar a superação de todas as falhas apontadas nos estudos da área. O uso interativo da tela só trará benefícios de ensino quando o aluno buscar realmente entender a relação entre as variáveis. O gráfico é importante, pois ajuda a visualizar conceitos que se tornariam muito abstratos se apenas falados. Porém, também pode criar a falsa ideia de que, apenas olhando a tela e mudando de parâmetro, será possível entender o assunto; por isso, cada elemento visual precisa ter um sentido claro e a lógica do modelo deve ficar acessível ao estudante.

Na proposta, o uso de dados sintéticos na aula tem o propósito de isolar as variáveis comparadas e mostrar de maneira clara o efeito da mudança de um único valor, mesmo que esse tipo de dado não sirva como previsão confiável sobre a realidade das pessoas. Já os dados coletados pela turma no segundo encontro mostram a origem dos números e permitem identificar a mudança ocorrida e qual parte do modelo produziu aquele efeito. O uso conjunto de tabelas, gráficos, fórmulas, textos curtos e controles, segundo \citeonline{ainsworth2006deft}, é uma forma de apresentar múltiplas representações que se complementam. No entanto, para que se tornem mais significativas, exigem atividades de interpretação e integração dos estudantes para responder à pergunta da aula, e não apenas juntar itens na mesma tela.

A visualização não dispensa o raciocínio do estudante. Para exemplificar, \citeonline{vartiainen2021machine} uniram a exploração de sistemas baseados em dados com a criação de ideias e explicações pelos alunos sobre o aprendizado de máquina; por isso, para cada controle da interface será realizada uma pergunta de previsão e explicação, evitando que a mudança vista no gráfico seja tomada como compreensão suficiente. Outro ponto observado é que qualquer ajuste realizado na tela (como mudança da inclinação da reta, posição da fronteira, vizinhos, condições ou centroides) requer que os alunos prevejam e interpretem a mudança, estabelecendo conexões entre os dados, a regra aplicada e o resultado. O papel visual é fundamental nos cinco modelos implementados para a visualização das informações aos alunos, pois cada modelo possui uma forma diferente de representação gráfica que vincula os dados e a regra à saída gerada. Para a regressão linear, por exemplo, os alunos podem visualizar os pares da tabela no plano cartesiano e a reta de tendência ou previsão. Já para o modelo de \textit{SVM}, os alunos podem ver a fronteira que separa as classes e a margem dos pontos mais próximos da divisão. De maneira similar, no modelo \textit{k}-NN, os alunos podem visualizar a posição do ponto consultado para medir a distância e acompanhar a votação dos vizinhos mais próximos. Os alunos também podem verificar o caminho das condições testadas nas árvores de decisão e os centroides móveis que reorganizam os grupos formados pelo modelo \textit{k}-means, longe de serem apenas pontos decorativos.

Para cada módulo, a atividade adapta a estratégia proposta por \citeonline{white1992probing} de prever, observar e explicar: o professor orienta os alunos a anotar uma hipótese oral ou escrita antes de iniciar o uso do controle para, posteriormente, checar no gráfico a medida calculada e a decisão tomada, explicando se o resultado confirmou a expectativa. Esse ciclo consiste em evitar a tentativa e erro aleatória em que parâmetros são mexidos sem relação clara entre entrada e saída. O ato de prever é importante porque abre a possibilidade de tornar visíveis ideias implícitas dos alunos. Por exemplo, se um aluno prevê que a inclusão de mais vizinhos no \textit{k}-NN não mudará a classe, mas o resultado da votação é diferente, o professor pode perguntar quais pontos passaram a ser considerados e por que a votação mudou. Esse é um ponto importante para que os alunos comecem a discutir e refletir sobre as características dos pontos de dados. No caso do algoritmo de agrupamento \textit{k}-means, por exemplo, o aluno pode prever que mover um centroide não fará diferença e ver os grupos mudarem na tela; nesse processo, a explicação retoma a ideia de proximidade entre os pontos. Errar a hipótese inicial também traz valor pedagógico, pois mesmo que o estudante não acerte o palpite inicial, ainda aprende com a análise da mudança.

As aulas estão previstas para acontecer no laboratório de informática do colégio, em que há conexão com a internet, dependendo da confirmação das condições e da estrutura disponível no local. A ideia é que cada aluno utilize um computador da sala ou um dos Chromebooks disponíveis, conforme a escola tiver, para realizar todas as tarefas durante o horário dos encontros. Dessa forma, evita-se que a falta de computador ou internet em casa possa prejudicar alguém. Para ajudar os estudantes que têm pouca prática ou receio em usar o computador, os botões e ajustes da tela deverão ser bem claros e fáceis de ler. Dessa forma, os alunos poderão saber a cada etapa da tarefa o que podem alterar. Já para os alunos que têm facilidade em mexer no computador, o foco da tarefa mudará: em vez de se preocuparem em operar a tela rapidamente, precisarão justificar conceitualmente os resultados, e o acompanhamento do professor, dando apoio individual quando necessário, ajuda a evitar que o manuseio técnico vire um obstáculo ou tome o lugar da discussão dos conceitos.

O controle interativo não se trata de um enfeite ou recompensa após uma explicação pronta, mas sim de uma ferramenta importante para obtenção dos resultados, permitindo ao estudante comparar os efeitos da mudança de um dado ou parâmetro em relação aos outros mantidos fixos. Por isso, durante a aula guiada com o professor, a orientação inicial é que se modifique apenas uma variável de cada vez, pois mexer em várias juntas dificilmente permitirá ao estudante saber qual foi a causa da mudança observada no gráfico. Após a aula guiada, os alunos terão a exploração livre para ampliar as comparações na tela. Nesse período, é importante que o professor retome a finalidade do experimento e solicite que os alunos expliquem o porquê de determinado resultado, e não apenas descrevam os elementos que estão aparecendo visualmente no gráfico.

A proposta apresenta o ensino das representações matemáticas em camadas graduais, partindo de uma situação concreta e que envolve a leitura de tabelas e gráficos simples, passando por cálculos básicos para posteriormente chegar à fórmula que resume a regra discutida. Dessa forma, tanto a equação da reta e a posição da fronteira quanto o cálculo de distâncias, as condições ou a média dos centroides entram para explicar mudanças visíveis na tela, evitando a cobrança de memorização solta. Trabalhar a Matemática em camadas graduais permite manter a precisão conceitual, definindo o que o estudante precisa para a compreensão inicial e o momento em que cada cálculo faz sentido para o iniciante. Também se destaca a importância de os alunos vivenciarem essas experiências antes de chegar à fórmula: é importante que entendam por que a proximidade entre os pontos importa no \textit{k}-NN antes de calcular a distância entre dois pontos, e que aprendam a comparar a previsão com o valor observado antes de discutir o erro de uma reta. Introduzida depois dessa análise visual, a fórmula ganha sentido prático porque nomeia uma relação que já foi examinada no gráfico durante o experimento.

\section{Avaliação da compreensão durante a intervenção}

Como já citado, abrir a discussão sobre a aprendizagem de \textit{machine learning} na Educação Básica é um campo em aberto nas pesquisas atuais. \citeonline{rauber2023assessing} realizaram um levantamento de instrumentos utilizados para avaliar a aprendizagem dos alunos, como testes e tarefas práticas, e verificaram a carência de formas de avaliação mais completas. Outro exemplo é o estudo de \citeonline{muhling2023novices} que realizou uma investigação sobre as ideias dos alunos em relação à inteligência artificial. Para tanto, os autores analisaram as respostas escritas abertas dos alunos antes e depois de uma oficina escolar. Mesmo com um design de pesquisa diferente, esta proposta também prioriza pedir que os estudantes explicitem de forma breve os conceitos empregados, para então analisar a compreensão desses conceitos, não se limitando a apenas contar as respostas certas e erradas.

\citeonline{gressevonwangenheim2022proposal} propuseram uma rubrica prática para medir os conceitos e habilidades inerentes ao \textit{machine learning} na escola, destacando a necessidade de transformar metas pedagógicas em critérios claros de observação. Mesmo criada para classificar fotos, sua rubrica é interessante, pois seus critérios podem ser adaptados aos cinco modelos apresentados neste trabalho. Ainda assim, saber de cor o nome de um algoritmo não prova que o estudante realmente entendeu a tarefa e como a regra se relaciona com os dados e o resultado. Outro ponto observado é que um acerto numérico pode ser alcançado de várias maneiras, como por exemplo, por meio do chute no gráfico, ao repetir as falas do professor, entre outras. Por isso, os instrumentos serão construídos para cobrar, de maneira breve, a justificativa do raciocínio realizado pelo estudante.

A compreensão dos alunos vai além de simplesmente marcar as opções certas nos testes de múltipla escolha. Por isso, o diagnóstico inicial e a parte equivalente no teste final exigirão que os alunos apresentem o mesmo nível de raciocínio e tamanho de tarefa. O diagnóstico no início servirá para mapear as ideias prévias da turma em relação aos conteúdos básicos de algoritmo, como por exemplo: acreditar que o computador é capaz de decidir sozinho sem precisar de dados; querer trocar a previsão de um valor por uma classificação; ou pensar que o agrupamento já vem com as classes prontas. O levantamento não tem a intenção de ranquear ou medir déficits. O instrumento final repetirá essas mesmas situações no primeiro bloco e acrescentará um segundo bloco com um problema novo para avaliar a transferência de um conceito.

De acordo com \citeonline{black1998assessment}, a avaliação formativa é vista como parte ativa do ensino, na qual as respostas dos alunos servem para que o professor possa verificar e ajustar seu caminho de explicações. Nesse sentido, as perguntas inseridas na ferramenta servem para dar um retorno rápido ao aluno e tentar puxar debates com a turma durante a aula, mas sem o sistema gravar nenhuma nota. Durante os testes na tela, o professor pode usar as hipóteses dos estudantes para verificar se existe alguma dúvida coletiva que precise de intervenção na hora (como, por exemplo, retomar o papel dos dados na margem da \textit{SVM} ou lembrar aos alunos da presença de rótulos no \textit{k}-NN e a ausência deles no \textit{k}-means). Após isso, o professor confronta esses registros com o diagnóstico e o instrumento final.

A análise dos dados se divide em quatro eixos principais: identificar a origem dos dados e a tarefa que foi pedida, ler as representações visuais, prever o efeito de ajustes e reconhecer as limitações das decisões automáticas. Para avaliar os alunos, a proposta adota uma rubrica com três níveis de resposta para organizar o material de análise entre: tentativas sem nenhuma relação válida, explicações parciais e respostas consistentes com o nível introdutório, de maneira a evitar reduzir todo o aprendizado a uma única nota numérica. Nas atividades práticas, os alunos precisarão, por exemplo, em uma situação sobre caminhada, separar os passos dados da distância percorrida para não tratar os números como dados genéricos. Também precisarão antecipar o que mudará caso movam os pontos, alterem a fronteira, troquem as condições ou mexam no número de vizinhos.

Como já citado, existe uma grande diferença entre dizer que o gráfico mudou porque o computador quis, e explicar de forma conceitual que aproximar a linha divisória diminuiu a margem dos exemplos mais próximos. Essa é uma das dimensões que serão consideradas na rubrica de avaliação e registradas em cada uma das quatro dimensões, acompanhadas de um comentário breve sobre o raciocínio do aluno para melhor orientação. O objetivo dessa análise qualitativa é acompanhar o processo para ver se os alunos estão entendendo os conceitos por trás da disciplina, sem precisar carimbar com rótulos definitivos de rendimento.

A apresentação dos resultados consistirá em mostrar resumos separados por dimensão e quadros de comparação entre o antes e o depois dos momentos equivalentes de cada dimensão, contando com contagens percentuais claras e pequenos trechos anônimos das falas dos alunos. A pesquisa acompanhará o que mudou e o que permaneceu no entendimento de cada aluno ao longo das aulas, sem comparar ou classificar os participantes entre si. A atividade de transferência inédita desenvolvida no final da intervenção e o questionário final de opinião dos estudantes serão discutidos em partes próprias, procurando manter a análise fiel às condições reais em que as aulas aconteceram na escola.

\section{Síntese do referencial e posicionamento da proposta}

A partir das diretrizes do currículo e dos artigos relacionados à Computação, percebe-se que ensinar \textit{machine learning} no colégio não se configura apenas em treinar a turma para usar programas prontos. Por isso, a proposta consiste em tratar o tema como um conteúdo de formação intrínseco à própria Computação. O plano de aula não tem o objetivo de adiantar o que é difícil na graduação, mas sim trabalhar as ideias básicas e os limites práticos que os alunos conseguem ver funcionando na tela. Para tanto, a ferramenta utilizada combina dados sintéticos com as tabelas criadas pela própria turma na sala, para que os estudantes possam mexer nos botões e verificar o resultado mudar na hora, bem como identificar de onde vieram aqueles números.

Como já citado no corpo do trabalho, durante a execução das tarefas, a tela pede ao aluno que faça uma previsão sobre o que vai acontecer antes de testar e ver o gráfico mudar, para em seguida explicar com suas palavras qual regra ele acredita que tenha produzido aquele resultado do computador. Na hora da avaliação do aprendizado dos alunos, a pesquisa não se prenderá apenas em contar quantos acertaram ou erraram nas atividades, nem busca provar regras gerais que possam servir para qualquer escola do país. O foco da análise será olhar com cuidado tudo o que cada um deles registrou durante as tarefas, procurando entender como suas ideias vão mudando dentro da realidade daquela sala de aula.

A literatura disponível já apresenta algumas propostas de ferramentas para ensino de IA que podem ser utilizadas nas escolas, porém muitas vezes são pouco detalhadas, não levam em consideração a realidade dos locais de aula e faltam avaliações mais rigorosas (\MakeUppercase{MARTINS; GRESSE VON WANGENHEIM}, \citeyear{martins2023findings}; \MakeUppercase{RAUBER; GRESSE VON WANGENHEIM}, \citeyear{rauber2023assessing}; \MakeUppercase{SANUSI} et al., \citeyear{sanusi2023systematic}). Por isso, este projeto tem a importância de contribuir com a área ao buscar investigar na prática escolar do próprio pesquisador uma sequência de aulas para uma turma do Ensino Médio, propondo a ligação da tela interativa com cinco algoritmos (regressão linear, \textit{SVM}, \textit{k}-NN, árvore de decisão e \textit{k}-means) a instrumentos com critérios claros para avaliação.

A qualidade científica de uma pesquisa está muito ligada ao alinhamento metodológico entre o referencial teórico escolhido e o produto educacional desenvolvido. Para tanto, é preciso que todos os elementos presentes na plataforma estejam bem planejados em função dos objetivos definidos nas aulas e nos instrumentos avaliativos, para que a proposta possa de fato contribuir com a realidade em que estará inserida. Outro ponto observado é que também é importante ter o cuidado metodológico de chegar às conclusões do trabalho respeitando exatamente o alcance dos dados obtidos com os estudantes.
